\chapter{1948:Arne Tiselius}

\label{chap:chemistry1948}

The Nobel Prize in Chemistry for the year 1948 was awarded to Arne Tiselius.

\textbf{Motivation:}

for his research on electrophoresis and adsorption analysis, especially for his discoveries concerning the complex nature of the serum proteins

\section{Arne Tiselius}

\subsection*{Early Life and Education}

Arne Tiselius was born on 1902-08-10 in Stockholm, Sweden.

\subsection*{Career and Major Research}

Tiselius studied chemistry at Uppsala University, where he received his doctorate under The Svedberg. He spent his career at Uppsala University. In 1937 he built a new apparatus for moving-boundary electrophoresis, in which proteins migrate in an electric field and can be measured as separate moving fronts.

\subsection*{Nobel Prize-winning Work}

The work for which Arne Tiselius was awarded the Nobel Prize in Chemistry in 1948 was described by the Nobel Committee as: "for his research on electrophoresis and adsorption analysis, especially for his discoveries concerning the complex nature of the serum proteins".

With his improved electrophoresis apparatus, Tiselius showed that the serum proteins form a complex mixture of several distinct components, notably albumins and different globulins. He also developed adsorption analysis, including methods for chromatographic separation of biological substances.

\subsection*{Significance and Impact}

Tiselius's electrophoretic separation of serum proteins revealed the complexity of the blood proteins and became a standard tool in biochemistry and clinical chemistry. Electrophoresis in its modern forms descends directly from his apparatus.

\subsection*{Later Career and Honors}

Tiselius continued at Uppsala University, where he also served as head of the Institute of Biochemistry. He later served as president of the Nobel Foundation. He died in 1971.

\subsection*{Legacy}

The electrophoretic methods introduced by Tiselius are fundamental to protein analysis, and his demonstration of the multiple components of serum proteins opened the field of clinical protein chemistry.

\subsection*{Modern Developments}

Tiselius's invention of electrophoresis made possible the separation and analysis of proteins, founding modern protein chemistry and clinical diagnostics. Electrophoresis remains a cornerstone of biochemistry and medicine, from hemoglobin and serum protein analysis to DNA sequencing and forensic genetics. His moving-boundary method evolved into today's gel, capillary, and microfluidic electrophoresis.

\subsection*{Selected Publications}

\begin{itemize}

\item \emph{A New Apparatus for Electrophoretic Analysis of Colloidal Mixtures}, Transactions of the Faraday Society, 1937.

\end{itemize}

\clearpage

\section*{Chapter Summary}

The 1948 prize recognized research on electrophoresis and adsorption analysis. Arne Tiselius's moving-boundary electrophoresis apparatus revealed that the serum proteins consist of several distinct components, and his methods became fundamental tools of protein chemistry.
